\documentclass[11pt]{amsart}
\usepackage{amsmath,amssymb,amsthm}
\usepackage[margin=1.15in]{geometry}
\usepackage{enumerate}
\usepackage{booktabs,array}
\usepackage[hidelinks]{hyperref}

\newtheorem{theorem}{Theorem}[section]
\newtheorem{lemma}[theorem]{Lemma}
\theoremstyle{remark}
\newtheorem{remark}[theorem]{Remark}

\newcommand{\doi}[1]{\href{https://doi.org/#1}{doi:\nolinkurl{#1}}}
\newcommand{\arxiv}[1]{\href{https://arxiv.org/abs/#1}{arXiv:\nolinkurl{#1}}}
\newcommand{\R}{\mathbb R}
\newcommand{\F}{\mathcal F}
\newcommand{\G}{\mathcal G}
\newcommand{\ceil}[1]{\lceil #1\rceil}
\newcommand{\floor}[1]{\lfloor #1\rfloor}
\setlength{\emergencystretch}{2.5em}

\begin{document}

\title{Hyperplane transversals of families with bounded packing number}
\author{Yitzchak Shmalo}
\date{25 November 2026}
\address{Einstein Institute of Mathematics, The Hebrew University of Jerusalem, Jerusalem, Israel}
\email{yspennstate@gmail.com}
\subjclass[2020]{52A35, 52C35, 05D15}
\keywords{hyperplane transversal, piercing number, $(p,2)$ property, packing number, Dol'nikov's theorem}

\begin{abstract}\label{abs:start}
McGinnis and Zerbib asked whether a finite family of compact convex sets in $\R^d$ with the $(p,2)$ property is pierced
by $p/d+O(1)$ hyperplanes. We show that a finite family of compact convex sets in
$\R^d$, $d\ge2$, whose largest subfamily of pairwise disjoint sets has $\nu$ members is pierced by
$\ceil{\nu/d}+2\ceil{\log_2 d}+3$ hyperplanes, and also by $\floor{\nu/d}+2d-2$ hyperplanes, which answers the
question in the affirmative when the constant is allowed to depend on $d$. Points on the moment curve show that the coefficient $1/d$ cannot be lowered. The additive term cannot be bounded
independently of $d$: hulls of the $s$-element subsets of points on the moment curve show that $\ceil{\log_2(\nu+1)}$
hyperplanes are sometimes needed, and this is the exact worst case when $d\ge 2^{\ceil{\log_2(\nu+1)}-1}$. The largest
excess over $\nu/d$, taken over all $\nu$, is $(1+o(1))\log_2 d$.
\end{abstract}

\maketitle

\section{Introduction}\label{sec:intro}

Let $\F$ be a finite family of nonempty compact convex sets in $\R^d$. Its \emph{packing number} $\nu(\F)$ is the
largest number of pairwise disjoint members of $\F$, and its \emph{hyperplane transversal number} $\tau(\F)$ is the
least number of affine hyperplanes whose union meets every member of $\F$. A family has the $(p,q)$ property if among
any $p$ of its members some $q$ have a common point; the $(p,2)$ property says exactly that $\nu(\F)\le p-1$. For
$d\ge2$ and $n\ge1$ let $T(d,n)$ be the largest value of $\tau(\F)$ over such families with $\nu(\F)\le n$.

Hyperplane transversals of families with a $(p,q)$ property go back to Alon and Kalai~\cite{AK95}, who proved that for
$p\ge q\ge d+1$ a family of compact convex sets in which some $q$ of every $p$ members have a common hyperplane
transversal is pierced by a number of hyperplanes that depends only on $p$, $q$ and $d$. For the $(p,2)$ property, projecting onto a line and
applying Gallai's theorem on intervals~\cite{Gal59} (see \cite[Theorem~1.6]{Zer24}) gives $T(d,n)\le n$. Zerbib~\cite[Theorem~1.2]{Zer24} used the KKM theorem
to prove that a finite family of compact convex sets in the plane with the $(p+1,2)$ property is pierced by
$\floor{p/2}+1$ lines, which is tight; this planar bound is a special case of a conjecture of McGinnis~\cite{McG24}. McGinnis and
Zerbib~\cite[Theorem~3.15]{MZ24} state the same bound $\floor{p/2}+1$ for hyperplanes in every dimension $d\ge2$, which
follows from the planar case by projecting onto a $2$-plane, since disjoint projections come from disjoint sets, and
ask~\cite[Question~3.16]{MZ24}:

\begin{quote}
Is it true that a finite family of compact convex sets in $\R^d$ with the $(p,2)$ property is pierced by $p/d+O(1)$
hyperplanes?
\end{quote}

Frick, Murray, Simon and Stemmler~\cite[Theorem~1.7]{FMSS24} proved that a family of polytopes in $\R^3$ with the
$(8,2)$ property is pierced by three planes. More generally they proved
that if no $2^k$ members of a finite family of polytopes in $\R^{2^{k-1}}$ are pairwise disjoint, then $k$ hyperplanes
pierce it; this is the case $m=1$ of~\cite[Theorem~1.5]{FMSS24}. Their methods are equivariant-topological and give the
right count when the number of hyperplanes is fixed and the dimension grows; they do not give a bound of the form
$\nu/d+O(1)$ in a fixed dimension.

If the implied constant is allowed to depend on $d$, Theorem~\ref{thm:upper} answers the question in the affirmative.
A constant independent of $d$ is impossible: by Theorem~\ref{thm:loss} the least admissible constant grows like
$\log_2 d$.

\begin{theorem}\label{thm:upper}
Let $d\ge2$, and let $\F$ be a finite family of nonempty compact convex sets in $\R^d$ with $\nu(\F)=\nu\ge1$. Then
\[
\tau(\F)\le \min\Bigl\{\Bigl\lceil\frac{\nu}{d}\Bigr\rceil+2\ceil{\log_2 d}+3,\ \Bigl\lfloor\frac{\nu}{d}\Bigr\rfloor+2d-2\Bigr\}.
\]
\end{theorem}

The second bound is the smaller one for $d\le5$ and the first for $d\ge7$; for $d=6$ the first is smaller when $6$
divides $\nu$, and the two agree otherwise. In dimension two Zerbib's bound is stronger. In
dimension three Theorem~\ref{thm:upper} improves on $\floor{\nu/2}+1$ for every $\nu\ge22$, and in each dimension from
$3$ to $10$ it does so for every $\nu\ge30$.

\begin{theorem}\label{thm:lower}
For all $d\ge2$ and $n\ge1$,
\[
T(d,n)\ \ge\ \max\Bigl\{\Bigl\lceil\frac nd\Bigr\rceil,\ \ceil{\log_2(n+1)}\Bigr\}.
\]
If $d\ge 2^{\ceil{\log_2(n+1)}-1}$, then $T(d,n)=\ceil{\log_2(n+1)}$.
\end{theorem}

For polytopes and $d=2^{k-1}$ the upper bound in the last sentence is the case $m=1$ of~\cite[Theorem~1.5]{FMSS24}
quoted above. Together, Theorems~\ref{thm:upper} and~\ref{thm:lower} give $T(d,n)/n\to1/d$ as $n\to\infty$ for each
fixed $d$, and show that no additive constant independent of $d$ is possible: $T(d,d)\ge\ceil{\log_2(d+1)}$, while
$d/d=1$. The size of the additive term is determined up to lower-order terms.

\begin{theorem}\label{thm:loss}
For $d\ge2$ let $A(d)=\sup_{n\ge1}\bigl(T(d,n)-n/d\bigr)$, and put $\ell=\ceil{\log_2 d}$ (fixed with this meaning
throughout) and $j=\ceil{\log_2(8\ell+16)}$. Then
\[
\ceil{\log_2(d+1)}-1\ \le\ A(d)\ \le\ \ell+2j+1,
\]
and in particular $A(d)/\log_2 d\to1$ as $d\to\infty$.
\end{theorem}

The proofs are elementary apart from one topological input. Two cutting operations are used. The first
(Lemma~\ref{lem:onecut}) splits one family, by one hyperplane, into two subfamilies with prescribed packing numbers
$a$ and $n-a-1$. The second (Lemma~\ref{lem:simul}) halves the packing numbers of up to $d$ families at once, provided
each of them is odd; it applies Dol'nikov's transversal theorem~\cite{Dol92} to the convex hulls of the large packings
of each family, the device by which Frick, Murray, Simon and Stemmler~\cite[Section~1]{FMSS24} relate transversals to
equipartitions. Halving repeatedly produces a binary tree of subfamilies. The point of Section~\ref{sec:schedule} is
that the nodes of such a tree can be scheduled breadth first, $d$ at a time, at a cost of only $\ceil{\log_2 d}$ extra
hyperplanes, and Section~\ref{sec:proof} pads an arbitrary packing number into such trees without losing a hyperplane
at each binary digit. For completeness, Section~\ref{sec:cuts} includes the projection proof of the transversal theorem
from the Borsuk--Ulam theorem~\cite{Bor33}.

The same bounds hold for compact connected sets, and for hyperplanes whose normals are restricted to a linear subspace
(Remark~\ref{rem:connected}). The paper is organized as follows; the symbols are listed in Appendix~\ref{app:notation}.

\begin{center}
\begin{tabular}{@{}lp{0.72\textwidth}@{}}
\toprule
Section & Contents \\
\midrule
\ref{sec:cuts} & lists and capacities; one cut (Lemma~\ref{lem:onecut}); Dol'nikov's theorem and simultaneous halving
(Lemma~\ref{lem:simul}) \\
\ref{sec:schedule} & breadth-first scheduling of a tree of halvings (Lemma~\ref{lem:schedule}) \\
\ref{sec:proof} & proof of Theorem~\ref{thm:upper} \\
\ref{sec:lower} & constructions on the moment curve; proof of Theorem~\ref{thm:lower} \\
\ref{sec:loss} & proof of Theorem~\ref{thm:loss} \\
\ref{sec:remarks} & connected sets, restricted normals, open questions \\
\bottomrule
\end{tabular}
\end{center}

\section{Two ways to cut a family}\label{sec:cuts}

A \emph{list} is a finite family of nonempty compact convex sets, possibly empty and possibly with repetitions; lists
are disjoint as lists even when their sets meet in space. A nonnegative integer $b$ is a \emph{capacity} for a list
$\G$ if $\nu(\G)\le b$ (with $\nu(\emptyset)=0$). If $b$ and $c$ are capacities for $\G$ and $\G'$, then $b+c$ is a
capacity for the concatenation of $\G$ and $\G'$, since a family of pairwise disjoint members splits into its parts in
$\G$ and in $\G'$. We say that a hyperplane $H$ \emph{processes} a list $\G$ into lists $\G^-$ and $\G^+$ if every
member of $\G$ that misses $H$ lies in exactly one of $\G^-$, $\G^+$; the members meeting $H$ are discarded.

\begin{lemma}\label{lem:onecut}
Let $n\ge1$ be a capacity for $\G$ and let $0\le a<n$. There is a hyperplane that processes $\G$ into lists with
capacities $a$ and $n-a-1$.
\end{lemma}

\begin{proof}
Fix a unit vector $u$. If $\G$ has no $a+1$ pairwise disjoint members, put every member missing an arbitrary
hyperplane in $\G^-$. Otherwise, for each family $P$ of $a+1$ pairwise disjoint members of $\G$ let
$m(P)=\max\{\langle u,x\rangle: x\in\bigcup P\}$, which exists by compactness, and let $t$ be the least of these finitely
many values, attained by $P_0$. Let $H=\{x:\langle u,x\rangle=t\}$. A member missing $H$ is contained in one of the two
open half-spaces, being convex; put those below $H$ in $\G^-$ and those above in $\G^+$. A family of $a+1$ pairwise
disjoint members of $\G^-$ would have $m(P)<t$, so $a$ is a capacity for $\G^-$. Every member of $\G^+$ lies above $H$,
and every member of $P_0$ lies in the closed half-space below $H$, so the members of $\G^+$ are disjoint from those of
$P_0$. Hence a family of pairwise disjoint members of $\G^+$ together with $P_0$ is pairwise disjoint in $\G$, and
$n-a-1$ is a capacity for $\G^+$.
\end{proof}

We use Dol'nikov's theorem in the following form; a family is \emph{intersecting} if any two of its members meet.

\begin{theorem}[Dol'nikov~\cite{Dol92}; see also {\cite[Theorem~1.2]{FMSS24}}]\label{thm:dol}
Let $\G_1,\dots,\G_k$, $k\le d$, be finite nonempty intersecting families of nonempty compact convex sets in $\R^d$. Then some
hyperplane meets every member of every $\G_i$.
\end{theorem}

\begin{proof}
For a unit vector $u$, the sets $\{\langle u,x\rangle:x\in C\}$, $C\in\G_i$, are pairwise intersecting closed
intervals, so they have the common interval $[\alpha_i(u),\beta_i(u)]$, where $\alpha_i$ is the largest left endpoint and
$\beta_i$ the smallest right endpoint. Both are continuous in $u$, since the families are finite and the support
function of a compact set is continuous. The midpoint $\mu_i=(\alpha_i+\beta_i)/2$ satisfies $\mu_i(-u)=-\mu_i(u)$.
The map $S^{d-1}\to\R^{d-1}$ whose first $k-1$ coordinates are $\mu_i(u)-\mu_1(u)$, $2\le i\le k$, and whose other
coordinates are zero is continuous and odd, so by the Borsuk--Ulam theorem~\cite{Bor33} it vanishes at some $u$. Then
$t=\mu_1(u)$ lies in every common interval, and $\{x:\langle u,x\rangle=t\}$ meets every member of every $\G_i$.
\end{proof}

\begin{lemma}\label{lem:simul}
Let $k\le d$, and let $2b_i+1$ be a capacity for the list $\G_i$, $1\le i\le k$. There is a hyperplane that processes
each $\G_i$ into two lists with capacity $b_i$ each.
\end{lemma}

\begin{proof}
Let $\mathcal K_i$ be the family of convex hulls $\operatorname{conv}\bigcup P$, where $P$ runs over the families of
$b_i+1$ pairwise disjoint members of $\G_i$. These hulls are compact and convex. Two of them meet: if
$\operatorname{conv}\bigcup P$ and $\operatorname{conv}\bigcup P'$ were disjoint, then so would be every member of $P$
and every member of $P'$, and $P\cup P'$ would be $2b_i+2$ pairwise disjoint members of $\G_i$. Apply
Theorem~\ref{thm:dol} to the nonempty families among $\mathcal K_1,\dots,\mathcal K_k$, and process each $\G_i$ by the
resulting hyperplane $H$ into the members below and above it. If $b_i+1$ pairwise disjoint members of $\G_i$ were on
the same open side of $H$, their hull would lie on that side and miss $H$. So $b_i$ is a capacity for both parts. (If
$\mathcal K_i$ is empty, $b_i$ is a capacity for $\G_i$ itself and there is nothing to prove.)
\end{proof}

For $b_i=0$, Lemma~\ref{lem:simul} gives a hyperplane meeting every member of $\G_i$: an intersecting list is pierced.

\section{Scheduling a binary tree}\label{sec:schedule}

A list with capacity $2^k-1$ is split by Lemma~\ref{lem:simul} into two lists with capacity $2^{k-1}-1$, and in
general into $2^r$ lists with capacity $2^{k-r}-1$ after $r$ rounds of splitting; at $r=k-1$ the lists have capacity
$1$, and the operations of the next level pierce them. We call one application of
Lemma~\ref{lem:simul} to one list an \emph{operation}; the operations form the \emph{perfect tree} of height $k$, with
$2^k-1$ nodes. We also use the \emph{truncated} tree, which stops before the last level: it has $2^{k-1}-1$ operations
and leaves $2^{k-1}$ lists with capacity $1$. A node is \emph{ready} once its parent has been performed.

By Lemma~\ref{lem:simul}, one hyperplane performs any $d$ ready operations at once, whatever their capacities and
wherever the lists lie in space. Members of other lists that the hyperplane happens to meet may be discarded or kept;
either way every capacity remains valid, since removing members of a list does not increase its packing number. Thus
the total count is governed by how many operations can be performed in parallel.

\begin{lemma}\label{lem:schedule}
A perfect or truncated tree with $s$ operations can be performed with at most $\ceil{s/d}+\ceil{\log_2 d}$ hyperplanes.
\end{lemma}

\begin{proof}
Number the nodes $1,2,\dots,s$ in breadth-first order, so that the parent of node $i\ge2$
is node $\floor{i/2}$. Perform the first $\ell$ levels one level per hyperplane, stopping if the tree ends; level $r$
has $2^{r-1}\le 2^{\ell-1}<d$ nodes, and these $\ell$ levels contain the nodes $1,\dots,2^\ell-1$. Split the remaining
nodes $2^\ell,\dots,s$ into consecutive batches of $d$ nodes, the last possibly smaller, and perform one batch per
hyperplane in order. If a batch starts at node $i$, then $i\ge2^\ell\ge d$, and the parent of its last node is at most
$\floor{(i+d-1)/2}<i$; so every node of the batch is ready when the batch is performed. The batches number at most
$\ceil{(s-2^\ell+1)/d}\le\ceil{s/d}$.
\end{proof}

\section{Proof of Theorem~\ref{thm:upper}}\label{sec:proof}

Write $n=\nu$. Several lists will be processed in sequence; at every stage every member of
$\F$ is either already met by a chosen hyperplane or lies in a current list, and every current list has a capacity.
When all current lists have capacity $0$ they are empty and $\F$ is pierced.

\subsection*{The first bound}
If $n=2^k-1$, Lemma~\ref{lem:schedule} applied to the perfect tree gives $\tau(\F)\le\ceil{n/d}+\ell$. Otherwise let
$2^k\le n<2^{k+1}-1$, so $k\ge1$, and put
\[
M=2^k-1,\qquad L=2^{k-1},\qquad a=n-2^k,
\]
so that $0\le a\le M-1$ and $n=M+a+1$. By Lemma~\ref{lem:onecut} one hyperplane processes $\F$ into a list with
capacity $M$ and a list with capacity $a$.

On the first list perform the truncated tree of height $k$: it has $M-L$ operations and leaves $L$ lists with capacity
$1$. Let $B$ be the least number of the form $2^r-1$, $r\ge0$, with $B\ge a$, and $t=B-a$. If $a=0$ then $B=t=0$.
Otherwise $(B-1)/2<a$, that is $a\ge(B+1)/2$, so $t\le(B-1)/2$; and $B\le M$ because $a<M$, so $t\le(M-1)/2=L-1$.
Concatenate $t$ of the $L$ lists of capacity $1$ with the list of capacity $a$; the result has capacity $a+t=B$, and we
perform its perfect tree, which has $B$ operations. The remaining $L-t$ lists of capacity $1$ are pierced $d$ at a time
by Lemma~\ref{lem:simul} with all $b_i=0$. By Lemma~\ref{lem:schedule} the number of hyperplanes is at most
\begin{align*}
1+\Bigl(\Bigl\lceil\frac{M-L}d\Bigr\rceil+\ell\Bigr)+\Bigl(\Bigl\lceil\frac Bd\Bigr\rceil+\ell\Bigr)+\Bigl\lceil\frac{L-t}d\Bigr\rceil
&\le 3+2\ell+\Bigl\lceil\frac{(M-L)+B+(L-t)}d\Bigr\rceil\\
&=3+2\ell+\Bigl\lceil\frac{n-1}d\Bigr\rceil,
\end{align*}
using $\ceil x+\ceil y+\ceil z\le\ceil{x+y+z}+2$ and $M+B-t=M+a=n-1$. This is at most $\ceil{n/d}+2\ell+3$.

\subsection*{The second bound}
Let $m=\floor{n/d}$. If $m=0$, then $n<d$, and $n$ applications of Lemma~\ref{lem:onecut} with $a=0$ reduce the
capacity of $\F$ to $0$; so $\tau(\F)\le n\le 2d-2$. Let $m\ge1$. Apply Lemma~\ref{lem:onecut} $d-1$ times, each time to
the current remainder, splitting off a list with capacity $m$; if the remainder already has capacity $m$ it is taken as
the next list, and empty lists fill the rest. Each split lowers the capacity of the remainder by $m+1$, and
$n-(d-1)(m+1)\le d(m+1)-1-(d-1)(m+1)=m$. So $d-1$ hyperplanes produce $d$ lists $\G_1,\dots,\G_d$, each with capacity
$m$.

If $m=2^k-1$, perform the perfect trees of all $\G_i$ simultaneously: one hyperplane performs the same node in all
$d$ trees by Lemma~\ref{lem:simul}. This uses $m$ hyperplanes. Otherwise apply the construction of the first bound to
each $\G_i$ with $m$ in place of $n$. Its quantities $M,L,a,B,t$ depend only on $m$, so after the $d$ separate cuts of
Lemma~\ref{lem:onecut} the remaining operations have the same shape in every $\G_i$: a truncated tree, a perfect tree
and $L-t$ piercings. Perform each operation in all $d$ copies with one hyperplane. The number of hyperplanes is
$(M-L)+B+(L-t)=m-1$. In all cases
\[
\tau(\F)\le (d-1)+d+(m-1)=\floor{n/d}+2d-2.
\]
This proves Theorem~\ref{thm:upper}.\qed

\section{Lower bounds}\label{sec:lower}

The moment curve $\gamma(\lambda)=(\lambda,\lambda^2,\dots,\lambda^d)$ meets every hyperplane in at most $d$ points, since
$\{\lambda:\langle u,\gamma(\lambda)\rangle=c\}$ is the zero set of a nonzero polynomial of degree at most $d$. Its points are in
convex position.

\begin{proof}[Proof of Theorem~\ref{thm:lower}]
For the first bound take $n$ distinct points on the moment curve as singletons. They are pairwise disjoint and a
hyperplane meets at most $d$ of them.

For the second let $k=\ceil{\log_2(n+1)}-1$, so that $2^k\le n$. Choose an integer $s>kd+1$ and points
$x_1,\dots,x_N$ on the moment curve with increasing parameters, where $N=(n+1)s-1$, and let $\F$ be the family of
convex hulls of all $s$-element subsets of $\{x_1,\dots,x_N\}$. The points are in convex position, so each hull meets
$\{x_1,\dots,x_N\}$ exactly in the $s$ points it is the hull of, and disjoint hulls come from disjoint sets of points; hence $\nu(\F)\le n$.
The hulls of the $n$ consecutive blocks $\{x_{(r-1)s+1},\dots,x_{rs}\}$, $1\le r\le n$, are pairwise disjoint because
their projections to the first coordinate are, so $\nu(\F)=n$. Suppose $k$ hyperplanes $H_1,\dots,H_k$ meet every
member of $\F$. At most $kd$ of the points lie on $\bigcup H_r$. The others are distributed among the at most $2^k$
sign classes $\{x:\operatorname{sign}(\langle u_r,x\rangle-c_r)=\epsilon_r,\ 1\le r\le k\}$, $\epsilon\in\{\pm1\}^k$,
each of which is convex. A sign class containing $s$ of the points would contain the hull of those $s$ points, which
then misses every $H_r$. So
\[
N\le kd+2^k(s-1)\le kd+n(s-1),
\]
whereas $N-n(s-1)=s+n-1>kd$. Hence $\tau(\F)\ge k+1=\ceil{\log_2(n+1)}$.

For the equality, let $h=\ceil{\log_2(n+1)}$ and suppose $d\ge2^{h-1}$. Then $2^h-1\ge n$ is a capacity for $\F$, and
the perfect tree of height $h$ has $2^{r-1}\le d$ nodes on level $r$. Performing one level per hyperplane by
Lemma~\ref{lem:simul} pierces $\F$ with $h$ hyperplanes.
\end{proof}

\section{The dependence on the dimension}\label{sec:loss}

\begin{proof}[Proof of Theorem~\ref{thm:loss}]
The lower bound is $A(d)\ge T(d,d)-1\ge\ceil{\log_2(d+1)}-1$ by Theorem~\ref{thm:lower}. For the upper bound let
$j=\ceil{\log_2(8\ell+16)}$ as in the theorem, and let $R$ be the least power of two with $R\ge d(2\ell+4)$. We call
a round \emph{full} if its hyperplane performs $d$ operations. Let $\F$ have $\nu(\F)=n$.

\emph{Large $n$.} Suppose $n\ge2R$. One hyperplane splits off, by Lemma~\ref{lem:onecut}, a list of capacity $2R-1$, leaving
a list of capacity $Q=n-2R$. The truncated tree of the first list has $R-1$ operations; by Lemma~\ref{lem:schedule} it
costs at most $(R-1)/d+\ell+1$ hyperplanes and leaves a reserve of $R$ lists of capacity $1$. Now run the construction
of Section~\ref{sec:proof} on the list of capacity $Q$: nothing if $Q=0$, one perfect tree if $Q=2^k-1$, and otherwise one
cut by Lemma~\ref{lem:onecut} followed by a truncated tree, a perfect tree and a set of piercings, with at most $Q$
operations in all. Scheduled as in Lemma~\ref{lem:schedule}, a tree with $s$ operations uses at most $\ceil{s/d}+\ell$
rounds and so leaves fewer than $d(\ell+1)$ idle positions in its rounds; the piercings leave fewer than $d$. In all,
fewer than $d(2\ell+3)\le R$ positions are idle, and we fill each of them with the piercing of a reserve list, which
Lemma~\ref{lem:simul} allows in the same round. Every round of this stage is then full except possibly one, made of the
reserve lists left over, so this stage uses at most $(Q+R)/d+1$ rounds, besides at most one cut. The total is at most
\[
1+\frac{R-1}d+\ell+1+\frac{Q+R}d+2=\frac{n-1}d+\ell+4 .
\]

\emph{Small $n$.} Suppose $n<2R$. Since $R<2d(2\ell+4)$, we have $n+1\le d(8\ell+16)\le2^{\ell+j}$. While the current
remainder has capacity $b>d$, split off by Lemma~\ref{lem:onecut} a list of capacity $2^r-1$, where $2^r\le b<2^{r+1}$;
the new remainder has capacity $b-2^r<b/2$. Since $n<d\,2^j$, there are at most $j$ such cuts. Round the final capacity
$b\le d$ up to the least $B=2^r-1\ge b$; then $B-b\le b\le d$. We now have a forest of perfect trees, with
$W\le n+d$ operations in all, since each cut removes one unit of capacity more than its tree has operations. Each tree
has height at most $\ceil{\log_2(n+1)}\le\ell+j$. In each round perform $d$ ready roots if there are at least $d$, and
all of them otherwise. A full round performs $d$ operations, and a round that is not full performs every current root,
so it lowers the largest remaining height by one; no round raises it. Hence there are at most $W/d+\ell+j$ rounds, and
\[
\tau(\F)\le j+\frac{n+d}d+\ell+j=\frac nd+\ell+2j+1 .
\]

Both cases give $T(d,n)-n/d\le\ell+2j+1$, since $j\ge5$. Finally $j=O(\log\log d)$ and $\ell/\log_2 d\to1$.
\end{proof}

\section{Remarks}\label{sec:remarks}

\begin{remark}\label{rem:connected}
The bounds of Theorems~\ref{thm:upper} and~\ref{thm:loss} hold for finite families of nonempty compact connected sets,
and also when the normals of the hyperplanes must lie in a given linear subspace $E$ of dimension $e\ge2$, with $e$ in
place of $d$. Let $\pi$ be the orthogonal projection onto $E$ and replace each set $C$ by $\hat C=\operatorname{conv}\pi(C)$,
a compact convex subset of $E$. For $u\in E$ we have $\langle u,\pi(x)\rangle=\langle u,x\rangle$, and $\langle u,C\rangle$
is an interval because $C$ is connected; so $\langle u,\hat C\rangle=\langle u,C\rangle$. Hence a hyperplane of $E$ with
normal $u$ meets $\hat C$ if and only if the hyperplane of $\R^d$ with the same equation meets $C$, and disjoint sets
$\hat C$ come from disjoint sets $C$, so $\nu$ does not increase. The moment curve in $E$ gives the lower bounds of
Theorem~\ref{thm:lower} with $e$ in place of $d$. For $e=1$ all hyperplanes are parallel, and the greedy choice of right
endpoints shows that $\nu$ of them suffice, and $\nu$ sets with pairwise disjoint projections show that $\nu$ are
needed.
\end{remark}

\begin{remark}
The problem of the best additive constant in a fixed dimension remains open. For $d=3$ the bounds give
$\max\{\ceil{n/3},\ceil{\log_2(n+1)}\}\le T(3,n)\le\floor{n/3}+4$ for $n\ge 3$, and~\cite[Theorem~1.7]{FMSS24} together with
Theorem~\ref{thm:lower} shows that the largest piercing number of a family of polytopes in $\R^3$ with packing number at
most $7$ is $3$. It is
also open whether $A(d)-\log_2 d$ is bounded.
\end{remark}

\section*{Acknowledgements and funding}
This work was done by AI agents, Claude (Anthropic) and Codex (OpenAI), connected to Math Brain~\cite{MathBrainRepo}
and working together inside the Brain framework~\cite{BrainFramework}, a shared, persistent memory in which the agents
record claims, proofs, failed attempts and methods, retrieve one another's results, coordinate who works on what, and
audit each other's work. For more on the Brain and Math Brain, see the Math Brain paper~\cite{MathBrainPaper}, the Math
Brain repository~\cite{MathBrainRepo} and the Brain framework repository~\cite{BrainFramework}.

The research presented in this paper was supported by the European Research Council (ERC) under the European Union's
Horizon Europe research and innovation programme (grant agreement No.~101041711), by the Simons Foundation as part of
the Collaboration on the Mathematical and Scientific Foundations of Deep Learning, by Heights Labs, by Convex Nexus
Capital, by the Israel Science Foundation (grant number 2258/19), and by the Israel Science Foundation (ISF Grant
4101/25).

\section*{Declaration of competing interest}
The author holds equity in Heights Labs and in Convex Nexus Capital and has a management role at Convex Nexus Capital.
These entities are included in the funding declaration above.

\section*{Data availability}
No data were generated or analysed in this work.

\section*{Declaration of generative AI and AI-assisted technologies}
This paper is part of a research project on AI proving problems. Codex agents (OpenAI) produced the constructions and
proofs, and Claude agents (Anthropic) checked the proofs and the references and wrote the manuscript. The question is
due to McGinnis and Zerbib~\cite{MZ24}, and its planar case was settled by Zerbib~\cite{Zer24}. The proofs rest on
Dol'nikov's transversal theorem~\cite{Dol92}, the Borsuk--Ulam theorem~\cite{Bor33} and the device of convex hulls of large
subsets from Frick, Murray, Simon and Stemmler~\cite{FMSS24}. The author directed the work and is responsible for the content of the paper and for decisions
concerning its dissemination.

\appendix
\section{Notation}\label{app:notation}

\begin{center}
\begin{tabular}{@{}>{\raggedright\arraybackslash}p{0.2\textwidth}>{\raggedright\arraybackslash}p{0.5\textwidth}>{\raggedright\arraybackslash}p{0.22\textwidth}@{}}
\toprule
Symbol or term & Meaning & First appears \\
\midrule
$(p,q)$ property & among any $p$ members some $q$ have a common point & abstract, page~\pageref{abs:start} \\
$\F$ & a finite family of nonempty compact convex sets in $\R^d$ & Section~\ref{sec:intro} \\
$\nu(\F)$ & packing number: the largest number of pairwise disjoint members of $\F$ & Section~\ref{sec:intro} \\
$\tau(\F)$ & the least number of hyperplanes whose union meets every member of $\F$ & Section~\ref{sec:intro} \\
$T(d,n)$ & the largest $\tau(\F)$ over families in $\R^d$ with $\nu(\F)\le n$ & Section~\ref{sec:intro} \\
$A(d)$ & $\sup_{n\ge1}\bigl(T(d,n)-n/d\bigr)$ & Theorem~\ref{thm:loss} \\
$\ell$ & $\ceil{\log_2 d}$ & Theorem~\ref{thm:loss} \\
$j$ & $\ceil{\log_2(8\ell+16)}$ & Theorem~\ref{thm:loss} \\
list & a finite family of nonempty compact convex sets, possibly empty & Section~\ref{sec:cuts} \\
capacity & an integer $b$ with $\nu(\G)\le b$ for the list $\G$ & Section~\ref{sec:cuts} \\
$H$ processes $\G$ & every member of the list $\G$ missing $H$ goes to exactly one of two lists; members meeting $H$ are discarded & Section~\ref{sec:cuts} \\
operation & one application of Lemma~\ref{lem:simul} to one list & Section~\ref{sec:schedule} \\
perfect, truncated tree & all operations that halve a list of capacity $2^k-1$ down to capacity $0$, or down to capacity $1$ & Section~\ref{sec:schedule} \\
ready & a node whose parent has been performed & Section~\ref{sec:schedule} \\
$\gamma$ & the moment curve $\lambda\mapsto(\lambda,\lambda^2,\dots,\lambda^d)$ & Section~\ref{sec:lower} \\
\bottomrule
\end{tabular}
\end{center}

The table lists the symbols and terms used in more than one section. Letters introduced inside a single proof or
construction or remark, such as $a$, $c$, $h$, $k$, $m$, $r$, $t$, $s$, $u$, $B$, $C$, $\hat C$, $E$, $L$, $M$, $N$, $P$,
$Q$, $R$, $W$, $\alpha_i$, $\beta_i$, $\mu_i$ and the lists $\mathcal K_i$, are defined where they appear and are local to it.

\begin{thebibliography}{99}
\bibitem{AK95} N.~Alon and G.~Kalai, \emph{Bounding the piercing number}, Discrete Comput. Geom. \textbf{13} (1995),
no.~3--4, 245--256, \doi{10.1007/BF02574042}.
\bibitem{Bor33} K.~Borsuk, \emph{Drei S\"atze \"uber die $n$-dimensionale euklidische Sph\"are}, Fund. Math.
\textbf{20} (1933), 177--190, \doi{10.4064/fm-20-1-177-190}.
\bibitem{Dol92} V.~L.~Dol'nikov, \emph{A generalization of the ham sandwich theorem}, Math. Notes \textbf{52} (1992),
no.~2, 771--779, \doi{10.1007/BF01236771}.
\bibitem{FMSS24} F.~Frick, S.~Murray, S.~Simon and L.~Stemmler, \emph{Transversal generalizations of hyperplane
equipartitions}, Int. Math. Res. Not. IMRN \textbf{2024} (2024), no.~7, 5586--5618, \doi{10.1093/imrn/rnad216};
\arxiv{2210.15423}, whose numbering is used here.
\bibitem{Gal59} T.~Gallai, \emph{\"Uber extreme Punkt- und Kantenmengen}, Ann. Univ. Sci. Budapest. E\"otv\"os Sect.
Math. \textbf{2} (1959), 133--138, Zbl 0094.36105.
\bibitem{McG24} D.~McGinnis, \emph{Piercing families of convex sets in the plane that avoid a certain subfamily with
lines}, Comput. Geom. \textbf{120} (2024), 102087, \doi{10.1016/j.comgeo.2024.102087}.
\bibitem{MZ24} D.~McGinnis and S.~Zerbib, \emph{Using the KKM theorem}, preprint, 2024, \arxiv{2408.03921}.
\bibitem{Zer24} S.~Zerbib, \emph{Bounds on piercing and line-piercing numbers in families of convex sets in the
plane}, Discrete Math. \textbf{347} (2024), no.~3, 113787, \doi{10.1016/j.disc.2023.113787}; \arxiv{2303.16240}, whose numbering is used here.
\bibitem{MathBrainRepo} Y.~Shmalo, \emph{Math Brain: public snapshot, code and tests}, 2026,
\url{https://github.com/yspennstate/math-brain-public/tree/5427734deee9be643d6dfbdd5e7d4baeeefe4a4a}.
\bibitem{MathBrainPaper} Y.~Shmalo, \emph{Math Brain: shared, auditable memory for multi-agent mathematical
research}, 2026, \doi{10.5281/zenodo.24150837},
\url{https://github.com/yspennstate/math-brain-public/tree/5427734deee9be643d6dfbdd5e7d4baeeefe4a4a/papers/math-brain-2026-10-24}.
\bibitem{BrainFramework} Y.~Shmalo, \emph{Brain Framework}, public software repository, 2026,
\url{https://github.com/yspennstate/brain-framework-public/tree/46b935d26d5a2893f7790b60466387dd8da47ef5}.
\end{thebibliography}

\end{document}
